\chapter{Introdução}

O Exame Nacional do Ensino Médio (ENEM) é a principal via de acesso ao ensino superior no Brasil, avaliando estudantes ao final da educação básica em quatro áreas do conhecimento: Linguagens e Códigos, Ciências Humanas, Ciências da Natureza e Matemática, além da redação. Os dados do ENEM 2024 revelam que, embora a média geral tenha subido de 543 para 546 pontos, a área de Matemática e suas Tecnologias caiu de 535 para 529 pontos, evidenciando uma dificuldade persistente dos estudantes nessa disciplina \cite{inep2025enem}.

A matriz de referência do ENEM exige que o estudante compreenda e utilize sistemas simbólicos para organização cognitiva da realidade, interpretação de dados e tomada de decisões \cite{matriz_enem}. Pesquisas indicam que as dificuldades não derivam apenas de lacunas de conhecimento, mas de esquemas cognitivos inadequados, baixa compreensão dos enunciados e desengajamento \cite{piton2018analise,soares2026diferencas}.

Diante disso, o Pensamento Computacional (PC) surge como habilidade analítica fundamental, descrita por \citeonline{wing2006} como tão indispensável quanto a leitura, a escrita e a aritmética, ao promover decomposição de problemas, reconhecimento de padrões, abstração e elaboração de algoritmos.

A Base Nacional Comum Curricular (BNCC) \cite{brasil2017} e a BNCC Computação \cite{brasil_bncc_computacao} estabelecem o PC como dimensão essencial para lidar com as tecnologias na sociedade contemporânea. Paralelamente, os elementos de gamificação -- pontos, insígnias e \textit{leaderboards} (PBL) -- podem elevar o engajamento comportamental e cognitivo dos estudantes \cite{silpasuwanchai2016developing}.

Este trabalho apresenta o \textbf{Cognitivo}, um software gamificado fundamentado nos quatro pilares do PC, projetado para aumentar o engajamento de estudantes em preparação para a prova de matemática do ENEM.